\section{Introduction}

Custody and event-log evidence matter because public trust is not only a count
arithmetic question. Voters and observers also ask when equipment was used,
which files were exported, whether media moved through expected hands, whether
seals and containers matched the records, and whether an observed process fits
the official explanation. Those are real election-administration questions, but
they are physical-process claims.

RCOUNT is not the right owner for those claims. A count-reconciliation verifier
can check that declared totals, subtotals, manifests, cast vote records, audit
records, and package hashes reconcile under named rules. It cannot conclude
from that reconciliation that ballots were always in proper custody, that a
seal was never broken, that a tabulator log is complete, or that an operator
performed an action for a stated reason.

This paper therefore reserves two downstream layers. RLOG is reserved for
normalized event and status logs. RCHAIN is reserved for physical
chain-of-custody evidence. The reservation is not an implementation claim. It
is a boundary claim: RCOUNT should remain honest about count evidence, and the
future custody and log layers should wait for source artifacts before they
harden schemas.
